% IEEEtran 风格：摘要用命令形式、关键词用环境、页眉用 \markboth、定理环境自定义
\documentclass[conference]{IEEEtran}
\usepackage{amsmath,amsthm}

\newtheorem{theorem}{Theorem}
\newtheorem{lemma}[theorem]{Lemma}

\begin{document}

\title{An IEEE Style Manuscript}
\author{Alice Chan and Bob Lee}
\markboth{IEEE Transactions on Testing, Vol.~1}{Chan and Lee}

\abstract{This abstract is written with the command form, which the LaTeX
reader drops entirely unless it is extracted first. It must survive.}

\begin{IEEEkeywords}
latex to word, mathtype, equation object
\end{IEEEkeywords}

\section{Introduction}
\label{sec:intro}

See Section~\ref{sec:intro}. The governing relation is

\begin{equation}
\sigma = E \varepsilon
\label{eq:hooke}
\end{equation}

\begin{theorem}
\label{thm:first}
Every well formed statement is numbered.
\end{theorem}

\begin{lemma}
\label{lem:first}
A lemma sharing the theorem counter follows it.
\end{lemma}

\begin{thebibliography}{9}
\bibitem{k1} A.~Author, \emph{A Book}, 2020.
\bibitem{k2} B.~Writer, ``A paper,'' \emph{Journal}, 2021.
\end{thebibliography}

\end{document}
